Es ist

\mavergleichskettealign
{\vergleichskettealign
{ (N (\gamma(t)) \times F(t))'
}
{ =} { (N (\gamma(t)))' \times F(t) +  N (\gamma(t)) \times F'(t)
}
{ =} { \left(  D N   \right)_{ \gamma(t)} \left(   \gamma'(t)   \right)  \times F(t) +  N (\gamma(t)) \times F'(t)
}
{ } {
}
{ } {
}
}
{}
{}{.}
Es ist zu zeigen, dass dieser Vektor stets senkrecht auf dem Tangentialraum $T_{\gamma(t)}Y$ steht. Da $F'(t)$ senkrecht auf dem
Tangentialraum steht, ist er linear abhängig zum Normalenvektor $N (\gamma(t))$ und daher ist der rechte Summand gleich $0$ nach
[[K^3/Kreuzprodukt/Eigenschaften/Fakt|Fakt   (3)]].
Im linken Summanden ist $F(t)$ tangential nach Voraussetzung und $\left(  D N   \right)_{ \gamma(t)} \left(   \gamma'(t)   \right)$ ist tangential nach
[[Hyperfläche/Weingartenabbildung/Tangentialraum/Fakt|Fakt]].
Daher ist nach
[[K^3/Kreuzprodukt/Eigenschaften/Fakt|Fakt   (6)]]
der linke Summand senkrecht auf dem Tangentialraum und somit ist insgesamt das Vektorfeld parallel.

 [[Kategorie:Latexseite]]
